% ============================================================================
% SEC05.TEX - Section 12.5: Code Optimization
% ============================================================================

\section{Code Optimization}

\begin{learningoutcomes}
\item Mengoptimalkan code untuk performa
\item Menghindari operations mahal di Update()
\item Menggunakan caching untuk GetComponent
\item Mengoptimalkan loops dan algorithms
\end{learningoutcomes}

\subsection{Best Practices}

\begin{lstlisting}[language={[Sharp]C}]
// BAD: GetComponent setiap frame
void Update()
{
    Rigidbody rb = GetComponent<Rigidbody>();
    rb.AddForce(Vector3.up);
}

// GOOD: Cache component
private Rigidbody rb;
void Start()
{
    rb = GetComponent<Rigidbody>();
}
void Update()
{
    rb.AddForce(Vector3.up);
}

// BAD: String concatenation di Update
void Update()
{
    text.text = "Score: " + score;
}

// GOOD: Update hanya saat berubah
void UpdateScore(int newScore)
{
    score = newScore;
    text.text = "Score: " + score;
}
\end{lstlisting}

\subsection{Optimization Tips}

\begin{itemize}
\item Cache GetComponent() di Start/Awake
\item Hindari operations mahal di Update()
\item Gunakan object pooling untuk Instantiate/Destroy
\item Minimize string operations
\item Use appropriate data structures
\end{itemize}

\begin{catatan}
Code optimization penting untuk performa. Profiler membantu mengidentifikasi code yang perlu dioptimalkan.
\end{catatan}

\subsection{Ringkasan Section}

\begin{itemize}
\item Cache GetComponent() untuk menghindari overhead
\item Hindari operations mahal di Update()
\item Gunakan object pooling
\item Minimize string operations
\item Profiler membantu mengidentifikasi masalah
\end{itemize}
